\section{Caracterização do problema}

\subsection{Problema da administração pública}

O problema central abordado neste trabalho consiste na dificuldade de transformar grandes acervos de dados sobre discursos legislativos em bases efetivamente consultáveis por meio de perguntas em linguagem natural, com respostas úteis, contextualizadas e verificáveis. Embora os discursos legislativos sejam públicos, seu volume, diversidade temática e dispersão temporal tornam a recuperação de argumentos, posicionamentos, temas recorrentes e referências históricas uma tarefa complexa quando realizada por mecanismos tradicionais de busca (KAVČIČ et al., 2024).

Na prática, esse quadro impacta diferentes atores institucionais. Gabinetes parlamentares podem necessitar recuperar rapidamente posicionamentos anteriores sobre determinado tema; equipes técnicas precisam localizar subsídios para notas, relatórios e pareceres; pesquisadores e jornalistas demandam acesso qualificado a discursos para fins analíticos; e cidadãos interessados em transparência e controle social podem se beneficiar de instrumentos que tornem o acervo mais acessível e inteligível (BANDEIRA et al., 2016). Na ausência de uma camada semântica de recuperação, o acesso ao conteúdo tende a ser fragmentado, lento e altamente dependente de conhecimento prévio sobre a estrutura documental e os termos empregados nos discursos (GAO et al., 2023; MARTELLOTE VIOLA et al., 2025).

Desse modo, o chatbot legislativo com RAG proposto neste trabalho procura atuar precisamente sobre essa lacuna, acrescentando uma camada de mediação informacional que combina recuperação semântica, contextualização documental e geração textual baseada em evidências.

\subsection{Partes interessadas}

As partes interessadas no projeto abrangem, em primeiro lugar, gabinetes parlamentares, consultorias legislativas e assessorias técnicas, que demandam acesso rápido e qualificado a discursos e posicionamentos anteriores. Incluem-se, também, equipes de transparência, gestão da informação e tecnologia, para as quais são particularmente relevantes aspectos de governança, manutenção, rastreabilidade e evolução da solução.

Além disso, pesquisadores das áreas de ciência política, direito, administração pública e ciência da informação constituem público potencial do sistema, assim como cidadãos e organizações da sociedade civil interessados em controle social e acompanhamento da atividade parlamentar (MARTELLOTE VIOLA et al., 2025; SAID, 2025).

\subsection{Critérios de sucesso}

Os critérios de sucesso definidos para este trabalho incluem: a capacidade de indexar uma base representativa de discursos parlamentares; a recuperação de trechos relevantes para perguntas factuais e analíticas; a geração de respostas em português fundamentadas prioritariamente no contexto recuperado; a apresentação de referências e metadados associados aos discursos utilizados; a automatização das etapas de importação e avaliação; e a reprodutibilidade da solução em ambiente local ou containerizado.

Tais critérios dialogam com recomendações da literatura para sistemas RAG em produção, que enfatizam avaliação contínua, reprodutibilidade e rastreabilidade como propriedades desejáveis em soluções voltadas a informação institucionalmente relevante (GAO et al., 2023; SAID, 2025).
